\section{Task and Benchmark Construction}

\subsection{Task Definition}

Each instance consists of claim text $x_t$, an associated image $x_v$, a decision cutoff $\tau$, and an evidence collection $E_{\leq\tau}$. The output label is $y \in \{\text{supported},\text{refuted},\text{insufficient}\}$. Evidence published after $\tau$ is unavailable by definition. The independent unit for splitting and uncertainty analysis is the event, not the individual post or document.

\subsection{Candidate Acquisition}

The proposed corpus uses English ClaimReview records concerning United States events from 2020--2026. GDELT supplies reproducible discovery metadata and Common Crawl supplies archived page observations. The collection retains cases with an identifiable claim time, recoverable image or archive, unambiguous professional rating, and sufficient timestamp metadata. Fact-check domains and documents containing the final verdict are excluded from the retrieval pool.

\subsection{Label Policy}

Only explicit site ratings are normalized. True or correct maps to supported; false, fake, or miscaptioned maps to refuted; and unproven, unsubstantiated, or no evidence maps to insufficient evidence. Mixed, partially true, satire, opinion, and ambiguous ratings are excluded. The project does not introduce new human labels. Automatically extracted evidence links are treated as silver evidence.

\subsection{Temporal Partitioning and Release}

Training covers 2020--2023, validation uses 2024, and final testing uses 2025--2026. Event identifiers, near-duplicate claims, and identical perceptual image hashes are group constraints. The proposed release contains URLs, cryptographic and perceptual hashes, timestamps, normalized labels, graph annotations, queries, and split manifests without redistributing copyrighted content.

Before treating \datasetname{} as a benchmark contribution, the filtered corpus must contain at least 2,000 cases, 300 per label, 50 test events, and two independent evidence clusters for at least 70\% of retained claims. Otherwise, it will be reported only as a stress-test collection.

